\section{Conclusion}
\label{sec:conclusion}

This paper presented target-local fusion V1 for prior-assisted multi-UAV
OTFS-ISAC sensing. The outer planner assigns each target to the UAV nearest its
predicted position, while the existing communication-calibrated adaptive C2F
selector performs detector-aligned greedy replay. The formulation therefore
changes where reports meet while the detector-PD objective independently
chooses its observation set under fixed packet, blocklength, and resource caps.

In 1000 paired trials, V1 achieved $P_D=0.9762$ and worst-target
$P_D=0.967$ from 12.101 observations, with only 0.751 remote reports,
0.481~kbit, and 0.801~ms of serial reporting. Its paired gain over
exact-marginal greedy was 0.0351. Its detection rate was statistically tied
with Sensing-SINR, while using 83.7\% less payload. A controlled fusion-rule
ablation showed that target-local planning disperses targets across fusion UAVs
and exposes shorter, more reliable reporting edges to the inner selector. The
evidence therefore supports a combined architectural conclusion: target-related
fusion planning plus adaptive two-stage selection---rather than attributing the
full improvement to a new greedy criterion.

V1 is a stable nominal version, not a universal endpoint. Large prediction
errors and sparse deployments remain its principal boundaries. The present
rate-threshold scan is weak because local evidence dominates remote reporting.
Future work will add fusion capacity and sensing-acquisition costs, validate the
analytic DD proxy at waveform level, model obstruction and correlated reports,
and optimize link-dependent blocklengths instead of a fixed report price.
